{\scriptsize
\setlength{\tabcolsep}{3pt}
\begin{longtable}{>{\raggedright\arraybackslash}p{2.19cm}>{\raggedright\arraybackslash}p{2.19cm}>{\raggedright\arraybackslash}p{2.19cm}>{\raggedright\arraybackslash}p{2.19cm}>{\raggedright\arraybackslash}p{2.19cm}>{\raggedright\arraybackslash}p{2.19cm}>{\raggedright\arraybackslash}p{2.19cm}}
\caption{Résultats de référence avec le nombre de composantes correctement spécifié ($K=3$)}\label{tab:2_10}\\
\toprule
\textbf{Modèle} & \textbf{Actives} & \textbf{ARI} & \textbf{Silh.} & \textbf{Itér.} & \textbf{Temps (s)} & \textbf{Poids} \\
\midrule
\endfirsthead
\toprule
\textbf{Modèle} & \textbf{Actives} & \textbf{ARI} & \textbf{Silh.} & \textbf{Itér.} & \textbf{Temps (s)} & \textbf{Poids} \\
\midrule
\endhead
\bottomrule
\endfoot
EM classique K=3 & 3 & 1,000 & 0,740 & 72 & 0,103 & 0,332 ; 0,333 ; 0,335 \\
EM variationnel Kmax=3 & 3 & 0,990 & 0,739 & 3 & 0,164 & 0,333 ; 0,335 ; 0,331 \\
\end{longtable}
}
